% analysis_compact.tex — §3 The Thinking Tax: Empirical Study (compact 9-page version)

All experiments use \textbf{Qwen3-8B}~\cite{qwen3_2025} on the full GSM8K test set ($n{=}1{,}319$), with cross-scale validation on Qwen3.5-9B and Qwen3.5-27B.
We control test-time compute via \texttt{max\_new\_tokens}, which caps \emph{total} output tokens for both modes equally.
Thinking mode's chain and answer share this budget; non-thinking mode uses it entirely for the answer.
Greedy decoding ($\tau{=}0$) is used throughout; we deliberately scope to the deterministic setting to isolate the truncation mechanism from sampling variance (the interaction between the thinking tax and stochastic decoding, including best-of-$K$ and self-consistency baselines, is an orthogonal question).
When thinking exhausts its budget without a final answer, a projection pass (up to 32 tokens) extracts a partial answer---\emph{favorable} to thinking mode (without it, think@512 drops from 65.2\% to 57.5\%; without \emph{any} answer extraction, to ${<}$6\%).

\paragraph{On fairness.}
One might object that capping thinking tokens is ``unfair.''
We note that (1)~in deployment, latency and cost budgets are real constraints~\citep{snell2024scaling}; (2)~we apply a projection pass that \emph{favors} thinking mode; and (3)~even a ``generous'' setup guaranteeing dedicated answer tokens trails nothink by 31.5\,pp (\S\ref{sec:crossover}).
The tax is not about answer truncation but about the reasoning chain consuming the \emph{entire} budget before reaching any answer.

\subsection{Finding 1: Non-Thinking Beats Thinking at All Matched Budgets}
\label{sec:finding:tax}

% Table 1: The Thinking Tax — Comprehensive comparison at matched budgets
\begin{table}[t]
\centering
\caption{\textbf{The Thinking Tax: Non-thinking mode dramatically outperforms thinking mode at matched token budgets.}
Results on GSM8K with Qwen3-8B.
Rows without a marker use the 200-sample subset (seed=42); starred ($\star$) rows use the full test set ($n{=}1{,}319$).
Qwen3.5-27B results use the full set.
\textbf{At every budget, non-thinking achieves higher accuracy while using fewer tokens.}}
\label{tab:thinking-tax}
\small
\begin{tabular}{ll rrrr}
\toprule
\textbf{Model} & \textbf{Config} & \textbf{Budget} & \textbf{Accuracy} & \textbf{Avg Tok} & \textbf{Early Stop} \\
\midrule
\multirow{10}{*}{\textbf{Qwen3-8B}}
& \texttt{nothink@128}      & 128 & 54.5\%  & 111 & 43.5\% \\
& \texttt{nothink@128$\star$} & 128 & 50.8\%  & 113 & --- \\
& \texttt{nothink@256}      & 256 & 89.0\%  & 140 & 92.0\% \\
& \texttt{nothink@256$\star$} & 256 & \textbf{87.5\%}  & \textbf{146} & 88.8\% \\
& \texttt{nothink@512}      & 512 & 94.0\%  & 145 & 99.5\% \\
\cmidrule{2-6}
& \texttt{think@128}        & 128 & 2.0\%   & 128 & 0.0\% \\
& \texttt{think@128$\star$}  & 128 & 3.0\%   & 128 & --- \\
& \texttt{think@256}        & 256 & 22.0\%  & 255 & 2.0\% \\
& \texttt{think@256$\star$}  & 256 & 18.0\%  & 255 & --- \\
& \texttt{think@512}        & 512 & 66.5\%  & 442 & 47.5\% \\
& \texttt{think@512$\star$}  & 512 & 65.2\%  & 460 & 37.4\% \\
\midrule
\multirow{3}{*}{\textbf{Qwen3.5-27B$\star$}}
& \texttt{think@128}        & 128 & 3.6\%   & 144 & 0.0\% \\
& \texttt{think@256}        & 256 & 7.9\%   & 272 & 0.0\% \\
& \texttt{think@512}        & 512 & 18.3\%  & 528 & 0.7\% \\
\bottomrule
\end{tabular}
\vspace{-2mm}
\end{table}


At every budget from 128 to 512, non-thinking mode uniformly dominates thinking mode (Table~\ref{tab:thinking-tax}).
\texttt{nothink@256} achieves \textbf{87.5\%} vs.\ only \textbf{18.0\%} for \texttt{think@256}---a \textbf{69.5\,pp} gap.
At $b{=}512$: 93.4\% vs.\ 65.2\% with projection ($+28.2$\,pp; 57.5\% without).
Non-thinking saturates by $b{=}512$ (93.4\%); thinking requires ${\sim}$2048 tokens for parity---a \textbf{4$\times$} budget multiplier.
The root cause is that thinking mode must fit its reasoning chain \emph{plus} its answer within $b$ tokens; at low budgets, the chain consumes all available tokens (token utilization analysis in Appendix~\ref{app:utilization}).
This does not imply CoT is useless---above the crossover (${\sim}$2048 on GSM8K), thinking recovers---but below this threshold, the reasoning overhead is a \emph{tax} that must be justified by problem difficulty.

\paragraph{Failure-mode decomposition.}
Table~\ref{tab:failure-decomp} decomposes the thinking tax into its constituent failure modes, providing causal evidence for the truncation mechanism.
At $b{=}256$, \textbf{98.6\%} of thinking responses are truncated before reaching a final answer, compared to only \textbf{11.2\%} in non-thinking mode.
Among the rare natural-stop thinking samples (1.4\%), accuracy is 100\%---the model \emph{can} reason correctly when given sufficient budget, but the chain format prevents it from finishing.
This rules out alternative explanations (prompt format, decoding mismatch): the only differentiating factor is whether the chain fits within the budget.
At $b{=}512$, natural-stop rate rises to 38.0\% with 99.0\% accuracy among completers; the truncated 62.0\% achieve only 32.2\%, dragging overall accuracy to 57.5\% (65.2\% with projection pass).
Under a stricter criterion ($<$95\% of budget, used in \S\ref{sec:theory}), the rate is 37.4\% with 96.3\% accuracy---the slight difference reflects near-boundary completions.

\begin{table}[t]
\centering
\caption{\textbf{Failure-mode decomposition} on GSM8K ($n{=}1{,}319$, Qwen3-8B).
Thinking mode's poor accuracy is almost entirely explained by truncation.
Think@512 overall is without projection; with projection it reaches 65.2\% (Table~\ref{tab:town-comparison}).
``Has final answer'' = truncated but a parseable answer was found in the output.
95\% bootstrap CIs from 10k paired resamples.}
\label{tab:failure-decomp}
\small
\begin{tabular}{@{}l@{\;\;}rr@{\;\;}rr@{}}
\toprule
& \multicolumn{2}{c}{\textbf{Budget 256}} & \multicolumn{2}{c}{\textbf{Budget 512}} \\
\cmidrule(lr){2-3} \cmidrule(lr){4-5}
\textbf{Metric} & \textbf{Think} & \textbf{Nothink} & \textbf{Think} & \textbf{Nothink} \\
\midrule
Natural stop rate & 1.4\% & 88.8\% & 38.0\% & 99.8\% \\
\;\;Acc.\ among natural stops & 100\% & 94.4\% & 99.0\% & 93.4\% \\
Truncated (hit budget) & 98.6\% & 11.2\% & 62.0\% & 0.2\% \\
\;\;Acc.\ among truncated & 16.8\% & 32.4\% & 32.2\% & --- \\
Has final answer & 0.4\% & --- & 5.3\% & --- \\
\midrule
\textbf{Overall accuracy} & 18.0\% & 87.5\% & 57.5\% & 93.4\% \\
\textbf{Tax (paired bootstrap)} & \multicolumn{2}{c}{69.5pp \ci{66.9}{72.0}} & \multicolumn{2}{c}{35.9pp \ci{33.2}{38.6}} \\
\bottomrule
\end{tabular}
\vspace{-2mm}
\end{table}

\subsection{Finding 2: Natural Stop as Confidence Oracle}
\label{sec:finding:natural-stop}

When thinking mode generates a complete chain and emits a final answer \emph{before} exhausting its budget, we say the sample exhibits a \emph{natural stop}.
At $b{=}512$ with Qwen3-8B, 37.4\% of samples terminate naturally; among these, accuracy is \textbf{96.3\%}---nearly perfect.
Among samples that exhaust the budget, accuracy is only 46.6\%, yielding a \textbf{49.7\,pp} gap (full analysis in Appendix~\ref{app:natural-stop-details}).
Unlike soft confidence scores, the natural-stop signal is \emph{binary}, \emph{endogenous}, and \emph{free}---it requires no logit access, hidden states, or calibration.
The effect generalizes across models: DeepSeek-R1-8B shows similar behavior (Appendix~\ref{app:deepseek}), with natural-stop rates increasing steadily with budget on MATH-500 (31.2\% $\to$ 76.2\% as $b$ scales from 1024 to 4096).
This signal forms the basis of \method{}'s routing mechanism (\S\ref{sec:method}).

\subsection{Finding 3: The Thinking Tax Scales with Model Size}
\label{sec:finding:model-size}

One might hope that the thinking tax diminishes with more capable models.
We test this by comparing Qwen3-8B, Qwen3.5-9B, and Qwen3.5-27B in thinking mode at identical budgets.
The results are counterintuitive.

% Table: Thinking Tax scales with model size (8B, 9B, 27B)
% Updated 2026-04-08 with gap-fill results: 8B nothink@512/1024/2048, 9B think@512/nothink@256/512, 8B think@1024
\begin{table}[t]
\centering
\caption{\textbf{The Thinking Tax worsens with model size} on GSM8K ($n{=}1{,}319$).
At budget=512, thinking-mode accuracy collapses from 56.9\% (8B) to 18.3\% (27B),
while non-thinking accuracy remains high across all scales:
8B=93.1\%, 9B=93.2\%, 27B=\textbf{95.5\%}.
The 9B model exhibits the largest tax (\textbf{77.7\,pp}), confirming that
inverse scaling is driven by chain-length growth, not capability loss.}
\label{tab:model-size-scaling}
\small
\begin{tabular}{l rr rr rr}
\toprule
 & \multicolumn{2}{c}{\textbf{8B}} & \multicolumn{2}{c}{\textbf{9B}} & \multicolumn{2}{c}{\textbf{27B}} \\
\cmidrule(lr){2-3} \cmidrule(lr){4-5} \cmidrule(lr){6-7}
\textbf{Budget} & \textbf{Think} & \textbf{NoThink} & \textbf{Think} & \textbf{NoThink} & \textbf{Think} & \textbf{NoThink} \\
\midrule
128 & 3.0\%  & 50.8\% & ---    & ---    & 3.6\%  & 9.9\% \\
256 & 18.0\% & 87.5\% & 3.4\%  & 61.2\% & 7.9\%  & 65.1\% \\
512 & 56.9\% & 93.1\% & 15.5\% & 93.2\% & 18.3\% & \textbf{95.5\%} \\
1024 & 86.1\% & 93.1\% & 41.8\% & ---$^\ddagger$ & ${\sim}$49\%$^\dagger$ & --- \\
2048 & ---$^\ddagger$ & 93.1\% & 66.8\% & ---    & ---    & --- \\
\midrule
\multicolumn{7}{l}{\textit{Thinking Tax at budget=512 (NoThink@512 $-$ Think@512):}} \\
\midrule
Tax & \multicolumn{6}{c}{8B: \textbf{36.2\,pp} (93.1\%$-$56.9\%); \quad 9B: \textbf{77.7\,pp} (93.2\%$-$15.5\%); \quad 27B: \textbf{77.2\,pp} (95.5\%$-$18.3\%)} \\
\bottomrule
\end{tabular}
\vspace{1mm}

\footnotesize
$^\dagger$Checkpoint at $n{=}200$; full-dataset experiment pending.\\
$^\dagger$Checkpoint at $n{=}200$; full-dataset experiment pending.
$^\ddagger$8B think@2048 in progress.\\
\textit{Root cause}: natural-stop rates at budget=512: 8B=37.4\%, 9B${\approx}$0.5\%, 27B=0.7\%.\\
At $b{=}1024$: 8B think=86.1\% vs nothink=93.1\% (gap 7.0\,pp); crossover not yet reached.\\
Nothink saturates: 8B nothink@512/1024/2048 = 93.1\%/93.1\%/93.1\% (avg ${\sim}$152 tokens).\\
\vspace{-3mm}
\end{table}


At $b{=}512$, thinking-mode accuracy \emph{collapses} with model size: 8B achieves 65.2\%, 9B drops to 11.5\%, and 27B reaches only 18.3\% (Table~\ref{tab:model-size-scaling}).
The thinking tax for 9B and 27B (relative to 8B nothink@512 = 93.4\%) is \textbf{2.9$\times$} and \textbf{2.7$\times$ larger} than for 8B, respectively.
The root cause is that larger models produce longer chains: at $b{=}512$, the 27B model's natural-stop rate is just \textbf{0.7\%} (vs.\ 8B's 37.4\%), so 99.3\% of its responses are truncated.
Crucially, the 27B model in \emph{non-thinking} mode achieves \textbf{95.5\%} at the same budget---the \emph{highest accuracy of any configuration}---confirming that model capability is intact and the tax is purely a format overhead.
Even quadrupling the budget does not close the gap: at $b{=}2048$, the 9B model reaches only 66.8\% in thinking mode---26.6\,pp below 8B \texttt{nothink@512} (93.4\%).
The tax worsens with scale, making adaptive allocation increasingly important for frontier models.

\paragraph{Summary.}
Our findings paint a coherent picture:
(1)~non-thinking dominates at all matched budgets;
(2)~natural stop predicts 96.3\% accuracy at zero cost;
(3)~the tax worsens with model scale.
The natural question is: can we \emph{think only when thinking helps}?
This is the goal of \method{}, introduced next.
